\chapter*{Abstract}
\addcontentsline{toc}{chapter}{Abstract}

Real-world training data can contain incorrect labels, noisy features, outliers and missing values. This project investigates how these corruptions, individually and in combination, affect classical machine learning models on tabular data.

A reproducible experimental framework evaluates Logistic Regression, RBF Support Vector Machine, Decision Tree and Random Forest on the Breast Cancer and Digits datasets. All fifteen non-empty combinations of four noise types are tested at five severity levels, using ten seeds and a stratified 70/30 split. The 6,080 fits apply noise only to training data and evaluate on a clean test set. Macro F1 is the primary metric and accuracy is secondary. A model's breaking point is the first level at which its mean Macro F1 falls to 90\% or less of its clean baseline. The Breaking Point Index and Robustness Score summarise performance, while a synergy analysis compares compound and individual noise effects.

In the official benchmark, label noise is the most damaging single corruption. Mean breaking levels fall from 5.281 for single noises to 3.750 for the four-way combination, counting a threshold not reached as level 6. Decision Tree is least robust on both datasets, while Random Forest has the highest mean Breaking Point Index and Logistic Regression retains the highest mean share of its clean Macro F1.

A separate demonstration uses all 581,012 records from UCI Covertype, a real-world forest-cover classification dataset. Across 16 Decision Tree fits using one seed, clean Macro F1 is 0.8950. The breaking point is level 1 for label noise and for label noise combined with missing values, and level 3 for missing values alone. A FastAPI service and React application support result exploration, a Covertype preset, and custom CSV uploads.
